\subsection{Twelf Wrapup}
Explaining all of the details of the Twelf code would
explode the length of this paper.
So in this section, we conclude the Twelf discussion by
just mentioning what each file contains and explaining
how Twelf helps with proving theory about languages.

The \code{sources.cfg} file tells Twelf the files to read
and the order in which to process them.
The \code{evaluation.elf} file contains the dynamic semantics
of $\minilang$.
\code{preservation.elf} contains the preservation theorem
and its proof.
\code{progress.elf} contains the progress theorem
and its proof.

Theorems are encoded in a similar manner to
(for all)-(there exists) queries.
For example, below is the encoding of the progress theorem.

\begin{verbatim}
%theorem  
progress :
  forall* {E} {T}
  forall {O:of E T} exists {NS:not_stuck E}
  true.
\end{verbatim}

The \code{forall*} line declares the variables in the theorem.
The next line states that for every derivation of the typing judgment,
\code{of E T}, where that derivation is bound to \code{O}, there
exists a derivation of the judgment that expression \code{E} is
not stuck (\code{NS:not\_stuck E}).
The following two rules define the \code{not\_stuck} judgment.

\code{not\_stuck/val: not\_stuck E <- value E.} states that if
\code{E} is value, then \code{E} is not stuck.

\code{not\_stuck/step: not\_stuck E <- step E E'.} states that if
there exists an \code{E'} such that \code{step E E'},
then \code{E} is not stuck as well.
The encoding of Twelf proofs are even more cryptic, so we skip its
discussion.

Finally, the benefit of using Twelf is that it can
check your proofs and occasionally derive proofs.
For example, the file \code{test\_typing.elf} provides two
example judgments that can be derived by Twelf.

\begin{verbatim}
%query * 1 D1 : of (estr (a , b , c , a , eps)) string.
%query * 1 D2 : of (add (enat (s z)) (enat z)) num.
\end{verbatim}

The first query ask Twelf to derive that the type of
the \code{str[abca]} is \code{str},
and save that derivation in \code{D1}.
The second query ask Twelf to derive that the
type of \code{+(num[1]; num[0])} is \code{num},
and save that derivation in \code{D2}.

Below is the output of Twelf from those queries.

\begin{verbatim}
loadFile test_typing.elf
[Opening file test_typing.elf]
%query * 1

of (estr (a , b , c , a , eps)) string.
---------- Solution 1 ----------
Empty Substitution.
D1 = of/str.
____________________________________________

%query * 1

of (add (enat (s z)) (enat z)) num.
---------- Solution 1 ----------
Empty Substitution.
D2 = of/add of/nat of/nat.
____________________________________________

[Closing file test_typing.elf]
%% OK %%
\end{verbatim}

Twelf finds derivations of both queries.
For the first query, it finds the derivation
\code{D1 = of/str.} saying to just apply
the \code{of/str} rule/function term.
\footnote{It is actually the case that
\code{D1 = of/str (a , b , c , a , eps).}
Twelf skips mentioning of the term
``\code{(a , b , c , a , eps)}'' because
it is passed as an implicit argument to
\code{of/str}.
Twelf only mentions explicit arguments in
its output.~\cite{Twelf-ConvertImplicit}
describes how to convert between implicit
and explicit arguments.}
For the second query,
it finds the derivation
\code{D2 = of/add of/nat of/nat.}
Derivation \code{D2} says:
\begin{enumerate}
\item Apply \code{of/nat} to \code{enat (s z)}
to return a derivation of \code{of (enat (s z)) num}.
\item Apply \code{of/nat} to \code{enat z}
to return a derivation of \code{of (enat z) num}.
\item Apply \code{of/add} to the derivations
of \code{of (enat (s z)) num} and \code{of (enat z) num}
to return a derivation of
\code{of (add (enat (s z)) (enat z)) num}.
\footnote{Again the expression terms are not
mentioned in the output because they are implicit
arguments.}
\end{enumerate}

Twelf could not derive the proof of the preservation theorem,
but it could let you know if your proof was incorrect.
For example, lines 13--17 in \code{preservation.elf}
prove preservation for the typing-evaluation rule combination
(T.4, D.1).
If I were to remove these lines, Twelf would return the following
error output:

\begin{verbatim}
preservation.elf:69.8-69.11 Error: 
Coverage error --- missing cases:
{E1:exp} {E2:exp} {E3:exp} {O1:of (add E1 E2) num} {S1:step E1 E3}
   {O2:of (add E3 E2) num}
   |- preservation O1 (step/add1 S1) O2.
\end{verbatim}

The error message lets us know that we forgot to prove preservation
for the case when we could derive \code{of (add E1 E3) num}
and \code{step E1 E3} ($e_1 \stepto e_3$); for this case,
we would need to produce a derivation of \code{of (add E3 E2) num}.

In summary, Twelf can aid with deriving simple judgments and checking
that our proofs of language properties are correct.
If our proofs are not correct, Twelf can let us know that are proofs
contain a mistake and provide an error message to help ``debug'' our proof.
